\section{Prime-supported rational gaps and two-row escape}
\label{sec:prime-support}

A second strand concerns the arithmetic conclusion of a rational linear-form
argument.  Let \(a,q,A,B\in\Z\) with \(q>0\), and
write
\[
  \Lambda_{A,B}(a/q)=B\frac aq-A .
\]
Here \(q\) is the denominator used to present the rational target, whereas
\(B\) is a coefficient of an integral approximation row.  They play different
roles.

\begin{theorem}[prime-supported two-row gap]
\label{res:prime-supported-two-selector-gap}
Let \(\ell\) be prime and let
\((A_1,B_1),(A_2,B_2)\in\Z^2\).  Suppose
\[
 \ell\nmid q,\qquad
 \ell\mid B_1,\quad \ell\mid B_2,\qquad
 \ell^2\nmid A_1B_2-A_2B_1 .
\]
Then
\[
 \left|\Lambda_{A_1,B_1}(a/q)\right|\ge\frac1q
 \quad\hbox{or}\quad
 \left|\Lambda_{A_2,B_2}(a/q)\right|\ge\frac1q .
\]
The constant \(1/q\) cannot be increased uniformly.
\end{theorem}

\begin{proof}
If \(\ell\) divided both \(A_1\) and \(A_2\), then \(\ell^2\) would divide
both \(A_1B_2\) and \(A_2B_1\), contrary to the determinant hypothesis.
Choose \(i\in\{1,2\}\) with \(\ell\nmid A_i\).  If
\(B_i a-A_iq=0\), primality and \(\ell\mid B_i\) would force
\(\ell\mid A_i\) or \(\ell\mid q\), again a contradiction.  Hence the
integer \(B_i a-A_iq\) is nonzero, and
\[
 \left|\Lambda_{A_i,B_i}(a/q)\right|
 =\frac{|B_i a-A_iq|}{q}\ge\frac1q .
\]
For sharpness take \(a=1\), \(q=2\), \(\ell=3\),
\((A_1,B_1)=(1,3)\), and \((A_2,B_2)=(2,3)\).  The determinant is \(-3\),
while both absolute remainders are \(1/2\).
\end{proof}

\begin{corollary}[rational two-row no-double-decay]
\label{res:rational-two-selector-no-double-decay}
Let \((A_{1,n},B_{1,n})\) and \((A_{2,n},B_{2,n})\) be sequences of integral
rows satisfying
\[
 A_{1,n}B_{2,n}-A_{2,n}B_{1,n}\ne0
 \qquad\hbox{for every }n .
\]
Then the two sequences
\[
 \Lambda_{A_{1,n},B_{1,n}}(a/q),
 \qquad
 \Lambda_{A_{2,n},B_{2,n}}(a/q)
\]
cannot both tend to zero.
\end{corollary}

\begin{proof}
At each \(n\), the nonzero exterior determinant prevents both rational linear
forms from vanishing.  Every nonzero integral linear form at \(a/q\) has
absolute value at least \(1/q\).  If both sequences tended to zero, then
eventually both absolute values would be strictly smaller than \(1/q\), a
contradiction.
\end{proof}

This corollary is an obstruction at rational targets, not an irrationality
theorem.  An application must still construct noncollinear rows whose two
analytic remainders tend to zero under the contrary rationality hypothesis.

\begin{proposition}[one-row prime-supported gap]
\label{res:prime-supported-one-row-gap}
Let \(\ell\) be prime, \(\ell\mid B\), \(\ell\nmid A\), and
\(\ell\nmid q\).  Then
\[
 \left|\Lambda_{A,B}(a/q)\right|\ge\frac1q .
\]
The constant is sharp.
\end{proposition}

\begin{proof}
The numerator \(Ba-Aq\) cannot vanish: otherwise \(\ell\mid Aq\), and
primality would contradict \(\ell\nmid A\) and \(\ell\nmid q\).  Its
absolute value is therefore at least one.  Equality is attained by
\(a=1,q=2,\ell=3,A=1,B=3\).
\end{proof}

\begin{corollary}[one-row prime-power-supported gap]
\label{res:prime-power-supported-one-row-gap}
Let \(\ell\) be prime, let \(r\ge1\), and suppose
\(\ell^r\mid B\), \(\ell\nmid A\), and \(\ell\nmid q\).  Then
\[
 \left|\Lambda_{A,B}(a/q)\right|\ge\frac1q .
\]
\end{corollary}

\begin{proof}
Since \(r\ge1\), the prime \(\ell\) divides \(B\), and
Proposition~\ref{res:prime-supported-one-row-gap} applies.
\end{proof}

The prime-power hypothesis records stronger denominator data but does not
improve the sharp \(1/q\) conclusion.

\begin{proposition}[real unit-determinant height--decay tradeoff]
\label{res:unit-determinant-height-decay}
Let \(F\in\R\), let \((A_1,B_1),(A_2,B_2)\in\R^2\), and let
\(u,v,w,z\in\R\) satisfy \(|uz-vw|=1\).  Define the recombined rows by
\[
 (A'_1,B'_1)=u(A_1,B_1)+v(A_2,B_2),\qquad
 (A'_2,B'_2)=w(A_1,B_1)+z(A_2,B_2).
\]
Suppose
\[
 |u|,|v|,|w|,|z|\le H
\]
and both recombined linear forms have absolute value at most
\(\varepsilon\ge0\).  Then
\[
 |A_1B_2-A_2B_1|
 \le 2H\varepsilon\bigl(|B_1|+|B_2|\bigr).
\]
\end{proposition}

\begin{proof}
Writing \(e_i=B'_iF-A'_i\), cancellation of the \(F\)-terms gives
\[
 A'_1B'_2-A'_2B'_1=e_2B'_1-e_1B'_2.
\]
Its absolute value is at most
\(\varepsilon(|B'_1|+|B'_2|)\).  The determinant of the recombined rows is
\((uz-vw)(A_1B_2-A_2B_1)\), so its absolute value equals that of the original
determinant.  Finally each new denominator coordinate is at most
\(H(|B_1|+|B_2|)\), giving the displayed bound.
\end{proof}

The transformation is real and has determinant of absolute value one;
integer-unimodular changes of basis are a special case.  Thus the proposition
quantifies a genuine coefficient-height cost: a change of basis cannot make
both remainders small without paying in the size of its coefficients.

\begin{proposition}[zero-denominator selector collision]
\label{res:zero-denominator-selector-collision}
Let \(N\ge1\) and \(w_i\in(\Z/N\Z)^2\) for \(i<k\).  If every second
coordinate of \(w_i\) is zero and
\[
 N<2^k,
\]
then two distinct binary selectors have the same vector sum.
\end{proposition}

\begin{proof}
Map a binary selector to the first coordinate of its sum.  There are \(2^k\)
selectors and only \(N\) possible values, so two distinct selectors collide.
Their second-coordinate sums are both zero by hypothesis.
\end{proof}

The threshold is \(N<2^k\), rather than the ambient two-coordinate condition
\(N^2<2^k\), because the denominator coordinate has already collapsed.

\paragraph{Exact Lean identities.}
The six source-current declarations corresponding to the results above are,
in theorem order:
\begin{flushleft}\small
\path{rational_twoSelector_gap_of_prime_tail_support}\\
\path{rational_twoSelector_not_both_tendsto_zero}\\
\path{rational_integerLinearForm_gap_of_prime_support}\\
\path{rational_integerLinearForm_gap_of_prime_power_support}\\
\path{twoSelector_unimodular_height_decay_tradeoff}\\
\path{zmod_binary_collision_of_zero_denominator_coordinates}
\end{flushleft}
The formal statements have the same hypotheses and conclusions as the
displayed results.

\paragraph{Exact remaining hypotheses.}
\label{bdry:prime-support-remaining-hypotheses}
Nothing in this section proves that the actual \(q\)-Ap\'ery rows contain a
suitable prime in their denominator coordinates cofinally, proves the
first-power determinant condition for a cofinal pair of rows, or supplies the
analytic remainder estimates needed to invoke
Corollary~\ref{res:rational-two-selector-no-double-decay}.  These are the exact
remaining hypotheses between the proved arithmetic consequences and an
irrationality proof at \(3/2\).
